\chapter*{Wykaz ważniejszych oznaczeń i skrótów}
\addcontentsline{toc}{chapter}{Wykaz ważniejszych oznaczeń i skrótów}

\begin{description}
    \item[UTH] \emph{Ultimate Texas Hold'em}: badana odmiana pokera kasynowego
    \item[RL] \emph{reinforcement learning}: uczenie przez wzmacnianie
    \item[DQN] \emph{Deep Q-Network}: algorytm uczenia przez wzmacnianie oparty na aproksymacji wartości akcji
    \item[REINFORCE] algorytm uczenia przez wzmacnianie z rodziny policy gradient
    \item[EV] \emph{expected value}: wartość oczekiwana, tu średni zysk netto na jedno rozdanie
    \item[CI] \emph{confidence interval}: przedział ufności
    \item[SD] \emph{standard deviation}: odchylenie standardowe
    \item[SE] \emph{standard error}: błąd standardowy
    \item[POMDP] \emph{partially observable Markov decision process}: częściowo obserwowalny proces decyzyjny Markowa
    \item[ReLU] \emph{rectified linear unit}: funkcja aktywacji sieci neuronowej
    \item[RAW] wariant reprezentacji stanu oparty wyłącznie na surowym kodowaniu kart
    \item[FEATURES] wariant reprezentacji stanu rozszerzony o ręcznie zaprojektowane cechy pokerowe
\end{description}

\bigskip

\begin{description}
    \item[$Q(s,a)$] funkcja wartości akcji $a$ w stanie $s$
    \item[$V^{\pi}(s)$] funkcja wartości stanu $s$ dla polityki $\pi$
    \item[$\gamma$] współczynnik dyskontowania nagród
    \item[$\varepsilon$] parametr strategii eksploracji epsilon-greedy
\end{description}
